\section{Arbitrary Generators Reduce to Partitions}\label{sec:partitions}
From now on $E(x)=(x,-x)$, $K_n=[0,1]^n$, and $D_n=\{t\one_n:0\le t\le1\}$. Admissible representatives have the \emph{global} exact slice $E^{-1}P=D_n$ and at most $M$ finite generators. Write
\begin{equation}\label{eq:minimax}
L_{\mathcal O}(M,n)=\inf_{\substack{P=\Hu(G),\ |G|\le M\\ E^{-1}P=D_n}}
\ \sup_{\substack{T\subseteq K_n\text{ finite}\\\varphi\in\mathcal O}}
\left(\min_{D_n\cup T}\varphi-\min_{J_T^{K_n}(P)}\varphi\right).
\end{equation}
We use $\Llip$ for $1$-Lipschitz observations and $\Laff$ for affine observations with $\|w\|_1\le1$. Real generator coordinates are allowed; the budget does not charge bit complexity.

For a partition $0=t_0<\cdots<t_q=1$, define
\begin{equation}\label{eq:partition-rep}
\begin{split}
P_{\boldsymbol t}&=\Hu\{e^0,e^1,c(t_{i-1},t_i):1\le i\le q\},\\
e^0&=E(0),\quad e^1=E(\one_n),\quad c(a,b)=(a\one_n,-b\one_n).
\end{split}
\end{equation}
Every point of this hull has constant positive and negative blocks, so its slice is contained in $D_n$. Conversely, for $a\le t\le b$,
\[
E(t\one_n)=\max\{c(a,b),e^0-t\one_{2n},e^1-(1-t)\one_{2n}\},
\]
which proves exactness.

\begin{theorem}[Contained partition representative]\label{thm:partition}
If $E(D_n)\subseteq P=\Hu\{g^1,\ldots,g^M\}$, then $M\ge3$ and $P$ contains a partition representative with at most $M-2$ cells. For exact $P$, replacing it by that representative weakly improves every contextual observation loss. Consequently the infima in~\eqref{eq:minimax} equal the infima over partitions with at most $M-2$ cells.
\end{theorem}
\begin{proof}
For $g=(u,v)$, define its anchor interval
$I_g=[0,1]\cap[\max_j u_j,\min_j(-v_j)]$.
The intervals cover $[0,1]$: a normalized representation of $E(t\one_n)$ has a coefficient-zero generator dominated by that point. Two distinct generators have no positive-length anchor. A generator attaining coordinate $-1$ at $e^0$ satisfies $v_1\ge0$ and $u_1\le v_1$, so its anchor ends at most at zero. A generator attaining coordinate $+1$ at $e^1$ satisfies $u_1\ge1$ and $v_1\le u_1-2$, so its anchor starts at least at one. The two inequalities on $u_1-v_1$ are incompatible for the same generator. The positive-length anchors alone still cover $[0,1]$, since their finite closed union contains the complement of finitely many points. There are at most $M-2$ of them.

Greedily choose an anchor reaching farthest right from the current endpoint. This yields a subordinate partition using at most one cell per chosen anchor. If $[a,b]\subseteq I_g$, then
\begin{equation}\label{eq:anchor-containment}
c(a,b)=\max\{g,e^0-b\one_{2n},e^1-(1-a)\one_{2n}\}\in P.
\end{equation}
The endpoint vectors belong to $P$ even if they are not displayed generators. This normalized identity proves containment. For exact inputs, both slices equal $D_n$ and joining a contained representative yields a contained result with the same concrete reference minimum. This proves the minimax reduction.
\end{proof}

Containment holds even for an inexact initial enclosure; comparing losses with different initial slices does not follow. The reduction changes the ambient set. It is not redundancy deletion preserving that set, nor a restriction of competitors to symmetric generators.

\subsection{Local cell costs and the exact Lipschitz law}
The homogeneous coordinate is covered for $P_{\boldsymbol t}$ exactly on the union of the anchor cubes $[t_{i-1},t_i]^n$. Outside that union there is no positive deficit. Inside one cube, write $u=\min_j x_j$, $v=\max_j x_j$. If $u<v$, the missing positive coordinates are precisely those with $x_j>u$, and the missing negative coordinates those with $x_j<v$. All constant-block generators can supply only extremal coordinates; the two endpoints supply all covered extrema. This description agrees with $P_0$ on each cube, including shared boundaries. Hence either observation class has partition loss equal to the maximum of its cell costs.

By~\eqref{eq:radius} and~\eqref{eq:signed-radii}, the pointwise Lipschitz deficit is
$\min\{(v-u)/2,u,1-v\}$.
For $n\ge2$, its maximum over $[a,b]^n$ is
\begin{equation}\label{eq:lipcell}
V(a,b)=\min\{(b-a)/2,b/3,(1-a)/3,1/4\}.
\end{equation}
Indeed, a deficit at least $\eta$ is feasible iff
$\max(a,\eta)+2\eta\le\min(b,1-\eta)$, giving these four bounds. Two coordinate values suffice for attainment.

\begin{theorem}[Exact Lipschitz generator law]\label{thm:lipschitz}
For $n\ge2$ and $M\ge3$,
\[
\Llip(M,n)=\frac1{2(M-1)}.
\]
An optimal partition with $q=M-2\ge2$ cells has boundary widths $3\delta$ and interior widths $2\delta$, where $\delta=1/[2(q+1)]$. For $q=1$ its cost is $1/4$. For $n=1$ the loss is zero.
\end{theorem}
\begin{proof}
The stated widths give the upper bound by~\eqref{eq:lipcell}. For a partition with $q\le M-2$ cells and maximum cost $\ell<1/4$, each cell $[a,b]$ has $b-a\le2\ell$, or $b\le3\ell$, or $a\ge1-3\ell$. Group the cells satisfying the second condition, then the remaining cells satisfying the third, then all other cells. The first two groups each cover length at most $3\ell$ when nonempty; each remaining cell covers at most $2\ell$. Thus the total length is at most $2(q+1)\ell$. Since it is one, $\ell\ge1/[2(M-1)]$. If $\ell\ge1/4$ the conclusion is immediate. Theorem~\ref{thm:partition} makes the bound exhaustive. Finally $D_1=K_1$, so there are no spurious states after clipping in dimension one.
\end{proof}

\subsection{A classical partition certificate}\label{sec:partition-certificate}
Let $W(a,b)$ be any nonnegative cell cost monotone under interval inclusion. Given two weak $q$-cell partitions $\boldsymbol s,\boldsymbol t$, some cell of $\boldsymbol t$ contains the corresponding cell of $\boldsymbol s$: take the first $i\ge1$ with $t_i\ge s_i$. Thus
\begin{equation}\label{eq:partition-certificate}
\min_i W(s_{i-1},s_i)\le
\inf_{\boldsymbol t}\max_i W(t_{i-1},t_i)\le
\max_i W(s_{i-1},s_i).
\end{equation}
An equal-cost partition is therefore globally optimal. The lower certificate must have exactly $q$ cells; a partition with fewer cells suffices only for the upper bound. Weak partitions allow zero widths and encode smaller budgets by repeated knots.

This interval-containment and balancing principle is classical~\cite[Lemma~2.12]{tilli2019}. We use sufficiency of equal costs, not uniqueness or necessity that every optimizer equalize. Non-strictly monotone costs may have plateaus. The specific work below is to compute the contextual cell cost and construct a feasible equal-cost partition.
